\documentclass[11pt,letterpaper]{article}
\usepackage[T1]{fontenc}
\usepackage[utf8]{inputenc}
\usepackage{lmodern}
\usepackage[margin=1in,headheight=14pt]{geometry}
\usepackage{amsmath,amssymb,amsthm,mathtools}
\usepackage{microtype}
\usepackage{booktabs,array,longtable}
\usepackage{xcolor}
\usepackage{fancyhdr}
\usepackage{enumitem}
\usepackage{listings}
\usepackage[hyphens]{url}
\usepackage[colorlinks=true,linkcolor=black,citecolor=black,urlcolor=black,
 pdfauthor={Research note prepared with ChatGPT},
 pdftitle={Leaves of Large Type in the Numerical-Semigroup Tree},
 pdfsubject={A counterexample to Conjecture 4.4 of Chappelon, Ramirez Alfonsin, and Stamate}]{hyperref}
\usepackage[nameinlink,noabbrev]{cleveref}
\pagestyle{fancy}
\fancyhf{}
\fancyhead[L]{\small\scshape Leaves of large type}
\fancyhead[R]{\small Research note \textnormal{\textbullet} September 2026}
\fancyfoot[C]{\thepage}
\renewcommand{\headrulewidth}{0.3pt}
\setlength{\parskip}{3pt}
\setlength{\parindent}{1.25em}
\linespread{1.025}
\setcounter{tocdepth}{2}
\numberwithin{equation}{section}
\newtheorem{theorem}{Theorem}[section]
\newtheorem{lemma}[theorem]{Lemma}
\newtheorem{proposition}[theorem]{Proposition}
\newtheorem{corollary}[theorem]{Corollary}
\theoremstyle{definition}
\newtheorem{definition}[theorem]{Definition}
\newtheorem{example}[theorem]{Example}
\theoremstyle{remark}
\newtheorem{remark}[theorem]{Remark}
\newcommand{\N}{\mathbb N_0}
\newcommand{\Z}{\mathbb Z}
\newcommand{\ints}[2]{[#1,#2]_{\Z}}
\newcommand{\tail}[1]{[#1,\infty)_{\Z}}
\DeclareMathOperator{\PF}{PF}
\DeclareMathOperator{\Ap}{Ap}
\DeclareMathOperator{\Max}{Max}
\newcommand{\conjb}{\beta}
\lstset{basicstyle=\small\ttfamily,breaklines=true,columns=fullflexible,
 keepspaces=true,showstringspaces=false,frame=single,rulecolor=\color{black!25},
 xleftmargin=0.5em,xrightmargin=0.5em,aboveskip=10pt,belowskip=10pt}
\title{\vspace{-0.7cm}{\LARGE\bfseries Leaves of Large Type in the\\[4pt]
 Numerical-Semigroup Tree}\\[9pt]
 {\large A genus-34 counterexample and a sharp asymptotic order}}
\author{A research investigation prepared for Vladimir Reshetnikov}
\date{20 September 2026}

\begin{document}
\maketitle
\thispagestyle{empty}
\begin{abstract}
We give an explicit counterexample to Conjecture 4.4 of Chappelon, Ram\'irez
Alfons\'in, and Stamate on the type of leaves of the numerical-semigroup tree.
The numerical semigroup
\[
 S=\{0,26,27,29,31,32\}\cup\ints{39}{50}\cup\tail{52}
\]
has genus $34$, type $17$, and no minimal generator above its Frobenius number
$51$; the conjectured bound is $16$. Its validity is proved directly, without
relying on a search program. We then turn restricted additive bases into
infinite families of leaves with explicitly determined pseudo-Frobenius sets.
One elementary family has genus $11q+1$ and type $6q-1$ and contradicts the
conjecture for every $q\ge3$.

More generally, let $L(g)$ be the maximum type of a leaf of genus $g$. A sparse
basis construction, a genus-padding operation, and a counting argument in
Ap\'ery sets prove
\[
 L(g)=\frac23g-\Theta(\sqrt g).
\]
In particular, the extremal type-to-genus ratio tends to $2/3$, not $1/2$.
Explicit two-sided bounds hold at every genus $g\ge7$. The archive includes
an exact-arithmetic verifier, independently reconstructed Ap\'ery certificates,
finite enumeration data, the random area-selection record, and optional search
code. The mathematical arguments are self-contained; the computational tests
are additional checks, not substitutes for the infinite proofs.
\end{abstract}

\vspace{6pt}
\noindent\textbf{Status.} This note disproves the stated conjecture and proves
the displayed asymptotic order. It does not determine the exact function
$L(g)$ or its sharp square-root coefficient. The retrieved version of the
source paper still states the conjecture; a targeted search did not locate a
later resolution. Priority and the proofs have not been independently refereed.

\newpage
\tableofcontents
\newpage

\section{The problem and the outcome}\label{sec:problem}

The area was selected before choosing the problem, using a pseudorandom
generator seeded from operating-system randomness. A uniform draw from 24
areas selected \emph{semigroup theory}. The seed and the full sampling frame
are recorded in \cref{sec:reproduce} and in the accompanying JSON file.

The target is Conjecture 4.4 of Chappelon, Ram\'irez Alfons\'in, and
Stamate~\cite{CRS}. For $g=4q+r\ge2$, with $0\le r<4$, put
\begin{equation}\label{eq:oldbound}
 \conjb(g)=
 \begin{cases}
  2q-1,&r=0,1,\\
  2q,&r=2,\\
  2q+1,&r=3.
 \end{cases}
\end{equation}
The conjecture asserts that every leaf $S$ of genus $g$ satisfies
$t(S)\le\conjb(g)$. Notice the strict inequality
\begin{equation}\label{eq:beta-half}
 \conjb(g)<\frac g2.
\end{equation}
Thus a leaf with $2t(S)\ge g(S)$ necessarily contradicts the conjecture.

The counterexample proved in \cref{sec:small} has $(g,t)=(34,17)$.
The source paper tabulates leaf counts through genus $33$~\cite[Table 5]{CRS},
so the example occurs immediately beyond that reported range. Our own
exhaustive enumeration stops at genus $20$; we do not claim an independent
exhaustive proof that genus $34$ is globally minimal.

The counterexample suggests an extremal problem that admits a substantially
stronger answer. In this note, define
\begin{equation}\label{eq:Ldef}
 L(g)=\max\{t(S):S\text{ is a leaf and }g(S)=g\}.
\end{equation}
The symbol denotes a \emph{maximum type}, not a number of leaves. The sets in
\eqref{eq:Ldef} are nonempty for every $g\ge7$, as our constructions show, and
they are finite. Our main quantitative result is
\begin{equation}\label{eq:summarybounds}
 \frac{2g}{3}-\sqrt{24g-87}-2
 \ \le\ L(g)\ \le\
 \left\lfloor\frac{6g+2-\sqrt{24g-23}}{9}\right\rfloor,
 \qquad g\ge7.
\end{equation}
Both errors from $2g/3$ have square-root order. The proof has two independent
parts: a construction for the lower bound and an Ap\'ery-set counting argument
for the upper bound. Neither uses an unproved classification of numerical
semigroups or an assumed optimality result for additive bases.

\section{Definitions and finite certificates}\label{sec:definitions}

We use $\N=\{0,1,2,\ldots\}$ and
$\ints{a}{b}=\{n\in\Z:a\le n\le b\}$. The notation
$A+B$ means $\{a+b:a\in A,b\in B\}$.

\begin{definition}
A \emph{numerical semigroup} is an additive submonoid $S\subseteq\N$ with
finite complement. Except when describing the root of the semigroup tree,
we assume $S\ne\N$. Its gap set, genus, Frobenius number, conductor, and
multiplicity are, respectively,
\[
 G(S)=\N\setminus S,\qquad g(S)=|G(S)|,\qquad F(S)=\max G(S),
\]
\[
 c(S)=F(S)+1,\qquad m(S)=\min(S\setminus\{0\}).
\]
A positive element is \emph{primitive} if it is not a sum of two positive
elements of $S$. The primitive elements form the minimal generating set;
its cardinality $e(S)$ is the embedding dimension.
\end{definition}

To see why the primitive elements generate, repeatedly split a nonprimitive
positive element into two smaller positive elements. This process terminates
because positive integers decrease. Conversely, a primitive element must
belong to every generating set. A primitive element cannot exceed $F+m$:
if $n\ge F+m+1$, then $n=m+(n-m)$ and $n-m\ge F+1$ belongs to $S$.

\begin{definition}
The \emph{numerical-semigroup tree} has root $\N$. The parent of a nonroot
semigroup $S$ is $S\cup\{F(S)\}$. Equivalently, its children are obtained by
deleting one primitive element larger than $F(S)$. A semigroup is a
\emph{leaf} when it has no children, or, equivalently, when every primitive
element is at most its Frobenius number.
\end{definition}

For completeness, adjoining $F$ preserves closure: every sum involving $F$
and a positive semigroup element is larger than $F$. Deleting a primitive
$x>F$ also preserves closure, since a sum of two remaining positive elements
cannot be $x$. The new largest gap is $x$, so these operations give exactly
the stated parent-child relation.

\begin{definition}
A \emph{pseudo-Frobenius number} is a gap $p\in G(S)$ such that
\[
 p+s\in S\qquad\text{for every }s\in S\setminus\{0\}.
\]
Write $\PF(S)$ for this set and $t(S)=|\PF(S)|$ for the \emph{type}.
\end{definition}

The largest gap $F$ always belongs to $\PF(S)$, so $t\ge1$. To verify that a
gap is pseudo-Frobenius, it suffices to test addition by minimal generators.
Indeed, after the first generator is added the sum belongs to $S$, and
further additions preserve membership. Alternatively, testing all positive
elements at most $F$ is finite and sufficient, because larger additions
already exceed the Frobenius number.

These observations give small certificates for statements that initially
quantify over an infinite semigroup. Closure only needs to be checked below
the conductor. The leaf condition only needs decompositions of the interval
$\ints{F+1}{F+m}$. A non-pseudo-Frobenius gap needs just one positive element
whose addition leaves another gap.

\section{A counterexample of genus 34}\label{sec:small}

\begin{theorem}\label{thm:small}
Let
\begin{equation}\label{eq:smallS}
 A=\{26,27,29,31,32\}\cup\ints{39}{50},\qquad
 S=\{0\}\cup A\cup\tail{52}.
\end{equation}
Then $S$ is a leaf with
\[
 (m,F,c,e,g,t)=(26,51,52,17,34,17).
\]
Its minimal generators are exactly $A$, and
\begin{align}\label{eq:smallPF}
 \PF(S)={}&\{13,14,15,16,17,18,21,23\}\nonumber\\
 &{}\cup\{28,30,33,34,35,36,37,38,51\}.
\end{align}
Consequently, $t(S)=17>\conjb(34)=16$.
\end{theorem}

\begin{proof}
Every positive element of $S$ is at least $26$. The sum of two positive
elements is therefore at least $52$ and belongs to $S$. This proves closure;
the full tail proves cofiniteness. The multiplicity is $26$, and the largest
gap is $51$.

There are five elements in $\{26,27,29,31,32\}$ and twelve in
$\ints{39}{50}$. Hence $A$ has $17$ elements and
\[
 g=51-|A|=34.
\]
Every element of $A$ is primitive, because it is smaller than $2m=52$.
We show that there are no others.

Put $D=\{0,1,3,5,6\}$. Direct addition gives
$D+D=\ints{0}{12}$. Therefore
\[
 (26+D)+(26+D)=\ints{52}{64}.
\]
Also $26+\ints{39}{50}=\ints{65}{76}$ and $77=27+50$.
Thus every element of $\ints{52}{77}$ is decomposable. For $n\ge78$,
\[
 n=26+(n-26),\qquad n-26\ge52,
\]
so every remaining tail element is decomposable too. The minimal generating
set is exactly $A$, all of whose elements are below $F=51$. Thus $S$ is a leaf.

The gap set is
\begin{equation}\label{eq:smallgaps}
 G(S)=\ints{1}{25}\cup\{28,30,33,34,35,36,37,38,51\}.
\end{equation}
Every gap in the second set is pseudo-Frobenius, since adding a positive
element gives at least $28+26>51$. For $13\le p\le18$, adding one of
$26,27,29,31,32$ gives an element of $\ints{39}{50}$, and adding an element
of $\ints{39}{50}$ gives at least $52$. For $p=21$, the five smaller-generator
sums are $47,48,50,52,53$; for $p=23$ they are $49,50,52,54,55$.
Thus these eight low gaps are pseudo-Frobenius as well.

The remaining seventeen gaps each fail the defining condition, as the
witnesses in \cref{tab:witness} show. This proves \eqref{eq:smallPF} and gives
$t=8+9=17$. Finally, $34=4\cdot8+2$ gives $\conjb(34)=16$.
\end{proof}

\begin{table}[htbp]
\centering
\begin{tabular}{rrr@{\qquad}rrr}
\toprule
Gap $p$ & Element $a$ & Gap $p+a$ & Gap $p$ & Element $a$ & Gap $p+a$\\
\midrule
1&27&28 &10&26&36\\
2&26&28 &11&26&37\\
3&27&30 &12&26&38\\
4&26&30 &19&32&51\\
5&29&34 &20&31&51\\
6&27&33 &22&29&51\\
7&26&33 &24&27&51\\
8&26&34 &25&26&51\\
9&26&35 & & & \\
\bottomrule
\end{tabular}
\caption{Certificates that the gaps not listed in \eqref{eq:smallPF} are not
pseudo-Frobenius. Every witness $a$ is a minimal generator.}\label{tab:witness}
\end{table}

This is a complete finite disproof of the conjecture. In particular, neither
the families below nor any asymptotic reasoning is required to establish the
contradiction. The purpose of the remaining sections is to explain its
structure and identify the correct scale of the extremal problem.

\section{Reducing depth-two leaves to finite sets}\label{sec:binary}

The counterexample has conductor exactly twice its multiplicity. At this
depth, additive closure is automatic, and the nontrivial conditions become
finite sumset and translation conditions.

\begin{proposition}\label{prop:binary}
Let $M\ge1$, let $m=M+1$, and let $B\subseteq\ints{0}{M}$ satisfy
$0\in B$ and $M\notin B$. Define
\begin{equation}\label{eq:SB}
 S_B=\{0\}\cup(m+B)\cup\tail{2m}.
\end{equation}
Then $S_B$ is a numerical semigroup with multiplicity $m$, conductor $2m$,
and Frobenius number $2M+1$. It is a leaf if and only if
\begin{equation}\label{eq:sumcondition}
 \ints{0}{M}\subseteq B+B.
\end{equation}
For
\begin{equation}\label{eq:Tdef}
 T(B)=\{p\in B\setminus\{0\}:(p+B)\cap\ints{0}{M}\subseteq B\},
\end{equation}
its invariants and pseudo-Frobenius set satisfy
\begin{align}
 g(S_B)&=2M+1-|B|,\label{eq:binaryg}\\
 \PF(S_B)&=T(B)\mathbin{\dot\cup}
    \bigl(m+(\ints{0}{M}\setminus B)\bigr),\label{eq:binaryPF}\\
 t(S_B)&=M+1-|B|+|T(B)|.\label{eq:binaryt}
\end{align}
When \eqref{eq:sumcondition} holds, the minimal generating set is exactly
$m+B$.
\end{proposition}

\begin{proof}
Every positive element is at least $m$, and every integer at least $2m$
belongs to $S_B$. Since $m+M=2m-1$ is absent, the asserted conductor and
Frobenius number are exact. All elements of $m+B$ are less than $2m$, hence
primitive.

It remains to examine possible primitive elements in
$\ints{2m}{3m-1}$. A decomposition there cannot use a summand at least $2m$,
because the other positive summand would bring the total to at least $3m$.
Consequently, $2m+j$ is decomposable exactly when
$j=b+b'$ for some $b,b'\in B$. This proves the leaf criterion. Above this
interval, $n=m+(n-m)$ is always a decomposition inside the tail.

The positive elements below the conductor are exactly the $|B|$ elements
of $m+B$, proving \eqref{eq:binaryg}. Every high gap
$m+j$, where $j\in\ints{0}{M}\setminus B$, is pseudo-Frobenius: adding
$m$ already reaches the tail. A low gap $p\in\ints{1}{M}$ can be
pseudo-Frobenius only if $p+m\in S_B$, that is, only if $p\in B$.
The remaining conditions are
\[
 p+(m+b)\in S_B\quad(b\in B).
\]
For $p+b>M$ they hold automatically, and for $p+b\le M$ they say exactly
$p+b\in B$. This is \eqref{eq:Tdef}. Low and high gaps are disjoint,
so taking cardinalities proves the type formula.
\end{proof}

The sumset condition explains the design problem. One wants a relatively
small set that covers an initial interval by two-term sums, together with a
long interval near the top that is stable under many small positive
translations. The former prevents tail generators; the latter creates many
low pseudo-Frobenius numbers. The two demands can be separated cleanly.

\section{Restricted additive bases produce leaves}\label{sec:basis}

\begin{definition}
A \emph{restricted additive basis of order two and maximum $d$} is a set
$D\subseteq\ints{0}{d}$ such that
\begin{equation}\label{eq:restricted}
 D+D=\ints{0}{2d}.
\end{equation}
We assume $d\ge1$ and write $k=|D|$.
\end{definition}

This is the restricted two-summand basis condition studied in the restricted
postage-stamp problem; see Kohonen~\cite{Kohonen} for that terminology and
computational setting. We need no extremal theorem from that literature.
Equation \eqref{eq:restricted} forces $0,1,d\in D$: representing $0$ and
$2d$ forces the endpoints, and representing $1$ forces $1$.

\begin{theorem}[Basis-to-leaf construction]\label{thm:basis}
Let $D$ satisfy \eqref{eq:restricted}, and let $\delta\ge0$ be an integer. Put
\begin{equation}\label{eq:construction}
 \begin{gathered}
 M=4d+1+\delta,\qquad m=M+1,\\
 H=\ints{2d+1}{M-1},\qquad B=D\cup H,\\
 S(D,\delta)=\{0\}\cup(m+B)\cup\tail{2m}.
 \end{gathered}
\end{equation}
Then $S(D,\delta)$ is a leaf, its minimal generators are precisely $m+B$, and
\begin{align}
 m&=4d+2+\delta,&F&=8d+3+2\delta,&c&=8d+4+2\delta,
 \label{eq:familymF}\\
 e&=k+2d+\delta,&g&=6d+3-k+\delta,&t&=4d+3-2k.
 \label{eq:familygt}
\end{align}
Its pseudo-Frobenius numbers are given explicitly by
\begin{align}\label{eq:familyPF}
 \PF(S(D,\delta))={}&
 \left(\ints{M-2d}{M-1}\setminus
       \{M-a:a\in D\setminus\{0\}\}\right)\nonumber\\
 &{}\mathbin{\dot\cup}\left(m+(\ints{0}{2d}\setminus D)\right)
   \mathbin{\dot\cup}\{2M+1\}.
\end{align}
In particular, increasing $\delta$ increases the genus by the same amount
and leaves the type unchanged.
\end{theorem}

\begin{proof}
The sets $D$ and $H$ are disjoint, $0\in B$, and $M\notin B$.
For the leaf condition, $D+D$ covers $\ints{0}{2d}$,
$0+H$ covers $\ints{2d+1}{M-1}$, and
\[
 M=1+(M-1)\in D+H.
\]
Thus $B+B$ covers $\ints{0}{M}$, so \cref{prop:binary} proves the leaf
property and identifies the generators.

The interval $H$ has $2d+\delta$ elements, whence
$|B|=k+2d+\delta$. This gives the formulas for $e$ and $g$ as well as
\eqref{eq:familymF}. For the type it remains to determine $T(B)$.

First, if $p\in D\setminus\{0\}$, then
\[
 d+1\le p+d\le2d.
\]
This interval is disjoint from $B$, so such $p$ cannot belong to $T(B)$.
Second, if $p\in H$ and $p\le M-2d-1$, then $M-p\in H$. The sum
$p+(M-p)=M$ is outside $B$, so this $p$ also fails the condition.

It remains to consider $p\in\ints{M-2d}{M-1}$. This entire interval lies
in $H$ because $M\ge4d+1$. For $h\in H$ we have $p+h>M$.
For $a\in D$, a sum $p+a\le M$ is already above $2d$, and hence it
belongs to $B$ unless it equals the single missing value $M$. Therefore
\begin{equation}\label{eq:Tfamily}
 T(B)=\ints{M-2d}{M-1}\setminus
       \{M-a:a\in D\setminus\{0\}\}.
\end{equation}
The interval has $2d$ elements, and the excluded $k-1$ values are distinct
and lie in it. Thus
\[
 |T(B)|=2d-k+1.
\]
The high missing indices are exactly $(\ints{0}{2d}\setminus D)\cup\{M\}$,
so there are $2d+2-k$ high pseudo-Frobenius numbers. Their sum with
$|T(B)|$ is $4d+3-2k$, and \eqref{eq:familyPF} follows from
\eqref{eq:binaryPF}.
\end{proof}

\begin{corollary}\label{cor:familydefect}
Every semigroup from \cref{thm:basis} satisfies
\begin{equation}\label{eq:defectfamily}
 F+t=2g,\qquad
 \frac{2g}{3}-t=\frac{4k-3+2\delta}{3}.
\end{equation}
At $\delta=0$ one also has $2t-g=2d+3-3k$.
\end{corollary}

\begin{proof}
Substitute \eqref{eq:familymF} and \eqref{eq:familygt}.
\end{proof}

The first identity says that these examples attain the general inequality
$F+t\le2g$, proved below. The second says that the departure from $2g/3$
is controlled by the basis size, not directly by its maximum. Sparse bases
are therefore the natural way to approach the limiting ratio.

\section{Two explicit infinite families}\label{sec:families}

\subsection{An odd-progression family containing the counterexample}

For an integer $q\ge1$, let
\begin{equation}\label{eq:oddD}
 d=2q,\qquad D_q=\{0,2q\}\cup\{1,3,5,\ldots,2q-1\}.
\end{equation}
Then $|D_q|=q+2$ and $D_q+D_q=\ints{0}{4q}$. Indeed, sums of two odd
members give all even integers from $2$ through $4q-2$. The even endpoints
come from $0+0$ and $2q+2q$. Adding $0$ or $2q$ to the odd members gives
all odd integers in the required range.

\begin{corollary}\label{cor:odd}
The semigroup $S_q=S(D_q,0)$ is a leaf with
\begin{equation}\label{eq:oddparams}
 m=8q+2,\qquad F=16q+3,\qquad e=5q+2,\qquad
 g=11q+1,\qquad t=6q-1.
\end{equation}
For every $q\ge3$ it contradicts \eqref{eq:oldbound}.
\end{corollary}

\begin{proof}
Apply \cref{thm:basis}. The excess over half the genus is
\[
 2t-g=2(6q-1)-(11q+1)=q-3.
\]
For $q\ge3$, this is nonnegative; now use \eqref{eq:beta-half}.
\end{proof}

At $q=3$, $D_q=\{0,1,3,5,6\}$, $M=25$, and
$H=\ints{13}{24}$, recovering \eqref{eq:smallS}. Thus the genus-34 example
is not isolated. Its successors increase the discrepancy without bound,
and $t/g$ tends to $6/11>1/2$.

\begin{table}[htbp]
\centering
\begin{tabular}{rrrrrr}
\toprule
$q$ & $g$ & $t$ & $\conjb(g)$ & $t-\conjb(g)$ & $2t-g$\\
\midrule
1&12&5&5&0&$-2$\\
2&23&11&11&0&$-1$\\
3&34&17&16&1&0\\
4&45&23&21&2&1\\
5&56&29&27&2&2\\
6&67&35&33&2&3\\
7&78&41&38&3&4\\
8&89&47&43&4&5\\
\bottomrule
\end{tabular}
\caption{The first odd-progression examples. The entries are direct evaluations
of proved formulas, also checked by the supplied program.}\label{tab:odd}
\end{table}

\subsection{A sparse square-grid family}

A different elementary basis improves the limiting ratio to $2/3$. For
$r\ge1$, put $d=r^2$ and define
\begin{equation}\label{eq:gridD}
 D_r=\ints{0}{r-1}\cup\{0,r,2r,\ldots,r^2\}
       \cup\ints{r^2-r+1}{r^2}.
\end{equation}
The set contains $3r-1$ elements. It is symmetric under $a\mapsto r^2-a$.
For $0\le n\le r^2$, Euclidean division gives
\[
 n=qr+s,\qquad 0\le s<r,
\]
with both $qr$ and $s$ in $D_r$. Reflecting this representation about $r^2$
represents every $n\in\ints{r^2}{2r^2}$ as well. Consequently,
$D_r+D_r=\ints{0}{2r^2}$.

\begin{corollary}\label{cor:grid}
For $S_r=S(D_r,0)$, write $g_r=g(S_r)$ and $t_r=t(S_r)$. Then
\begin{equation}\label{eq:gridgt}
 g_r=6r^2-3r+4,\qquad t_r=4r^2-6r+5.
\end{equation}
These are leaves, and
\begin{equation}\label{eq:gridratio}
 t_r=\frac23g_r-4r+\frac73,\qquad
 \lim_{r\to\infty}\frac{t_r}{g_r}=\frac23.
\end{equation}
They contradict the original conjecture for every $r\ge4$.
\end{corollary}

\begin{proof}
The formulas follow by substituting $d=r^2$ and $k=3r-1$ into
\cref{thm:basis}. Moreover,
\[
 2t_r-g_r=2r^2-9r+6,
\]
which is positive at $r=4$ and strictly increasing for integer $r\ge4$.
The contradiction again follows from \eqref{eq:beta-half}.
\end{proof}

\begin{table}[htbp]
\centering
\begin{tabular}{rrrrrr}
\toprule
$r$ & $|D_r|$ & $g_r$ & $t_r$ & $\conjb(g_r)$ & $t_r-\conjb(g_r)$\\
\midrule
1&2&7&3&3&0\\
2&5&22&9&10&$-1$\\
3&8&49&23&23&0\\
4&11&88&45&43&2\\
5&14&139&75&69&6\\
6&17&202&113&100&13\\
8&23&364&213&181&32\\
10&29&574&345&286&59\\
20&59&2344&1485&1171&314\\
30&89&5314&3425&2656&769\\
\bottomrule
\end{tabular}
\caption{Sparse-grid leaves. The quadratic growth of genus and type, versus the
linear size of the additive basis, yields the limiting ratio $2/3$.}
\label{tab:grid}
\end{table}

\subsection{Why square-root defects arise in this construction}

There are at most $k(k+1)/2$ distinct unordered two-term sums from a set of
size $k$, allowing equal summands. A restricted basis therefore satisfies
\[
 \frac{k(k+1)}2\ge2d+1,
 \qquad
 k\ge\frac{\sqrt{16d+9}-1}{2}.
\]
Thus its size is necessarily at least of square-root order in $d$.
The bases \eqref{eq:gridD} have size $3\sqrt d-1$, which is sufficient for
our asymptotic purpose even without any claim of optimality. By
\eqref{eq:defectfamily}, these bounds translate directly into square-root
defects from $2g/3$ along the unpadded family.

This counting observation applies only to this construction. The next
section gives a separate argument showing that a square-root defect is
unavoidable for \emph{every} leaf, whether or not it has conductor $2m$ or
comes from a restricted basis.

\section{Universal upper bounds for the type of a leaf}\label{sec:upper}

\subsection{The reflection inequality}

The inequality in the following lemma is standard; it also appears as
Proposition 3.1 of the target paper~\cite{CRS}. We include a direct proof so
that all later bounds can be checked from the definitions.

\begin{lemma}\label{lem:reflection}
Every nontrivial numerical semigroup satisfies
\begin{equation}\label{eq:reflection}
 F+t\le2g.
\end{equation}
\end{lemma}

\begin{proof}
Let $N=S\cap\ints{1}{F}$. Its cardinality is $F-g$. The sets $N$ and
$\PF(S)\setminus\{F\}$ are disjoint. Reflection $x\mapsto F-x$ sends
their union injectively into $G(S)\setminus\{F\}$.

For $x\in N$, membership $F-x\in S$ would imply $F\in S$ by closure.
For $x\in\PF(S)\setminus\{F\}$, membership of the positive number
$F-x$ in $S$ would contradict the pseudo-Frobenius condition. All reflected
values are positive and smaller than $F$. Comparing cardinalities gives
\[
 (F-g)+(t-1)\le g-1,
\]
which is equivalent to \eqref{eq:reflection}.
\end{proof}

Since $t\ge1$, this also gives $F\le2g-1$. Thus there are only finitely many
numerical semigroups of any fixed positive genus: their gap sets are subsets
of the finite interval $\ints{1}{2g-1}$. This justifies the finiteness asserted
in the definition of $L(g)$.

\subsection{Ap\'ery sets and nonmaximal elements}

\begin{definition}
For the multiplicity $m$, the \emph{Ap\'ery set} is
\[
 W=\Ap(S;m)=\{s\in S:s-m\notin S\}.
\]
It contains the least element of $S$ in each residue class modulo $m$, hence
has $m$ elements and contains $0$. On $W$ use the partial order
\[
 u\preceq_S v\quad\Longleftrightarrow\quad v-u\in S.
\]
\end{definition}

The numerically largest Ap\'ery element is $F+m$. Indeed, it belongs to $S$
while subtracting $m$ gives the gap $F$; conversely, $w-m\notin S$ prevents
an Ap\'ery element $w$ from exceeding $F+m$.

\begin{lemma}\label{lem:aperyPF}
The map $p\mapsto p+m$ is a bijection from $\PF(S)$ to the maximal elements
of $(W,\preceq_S)$. In particular, precisely $t$ elements of $W$ are maximal,
and none of them is $0$.
\end{lemma}

\begin{proof}
For $p\in\PF(S)$, $p+m\in S$ and $p\notin S$, so $p+m\in W$.
If $w\in W$ strictly dominates $p+m$, write $w=p+m+s$ with positive
$s\in S$. Then $w-m=p+s\in S$, a contradiction.

Conversely, let $w$ be maximal in $W$. It is nonzero because $0$ is below
every element of $W$ and $m\ge2$. Every nonzero Ap\'ery element exceeds
$m$, so $p=w-m$ is a positive gap. If $p+s\notin S$ for some positive
$s\in S$, then $w+s\in W$, contradicting maximality. Thus $p\in\PF(S)$.
\end{proof}

Write
\begin{equation}\label{eq:a}
 a=m-1-t
\end{equation}
for the number of nonzero, nonmaximal Ap\'ery elements.

\begin{lemma}\label{lem:leafbasic}
For every leaf, $a\ge1$ and $F\ge2m-1$. Consequently, with
$h=2g-3t$, one has
\begin{equation}\label{eq:hlinear}
 h\ge2a+1\ge3,
 \qquad
 t\le\left\lfloor\frac{2g-3}{3}\right\rfloor.
\end{equation}
\end{lemma}

\begin{proof}
The Ap\'ery element $F+m$ lies above $F$ and therefore is not primitive in
a leaf. Write $F+m=u+v$ with $u,v\in S\setminus\{0\}$.
Both $u$ and $v$ belong to $W$: if $u-m\in S$, for example, then
$F=(u-m)+v\in S$, a contradiction. Both are nonmaximal since
$u\prec_S F+m$ and $v\prec_S F+m$. They may coincide, but at least one
nonzero nonmaximal element exists. Hence $a\ge1$.

If $F<2m-1$, the number $2m-1$ belongs to $S$ and is smaller than the sum of
any two positive elements. It is therefore a primitive element above $F$,
contrary to the leaf condition. Thus $F\ge2m-1$.

Combining this with \cref{lem:reflection} and $m=t+a+1$ gives
\[
 2g-3t\ge F-2t\ge2m-1-2t=2a+1.
\]
The final bound follows from $a\ge1$.
\end{proof}

\subsection{Counting decompositions above the Frobenius number}

The preceding linear bound identifies the candidate leading constant $2/3$,
but it does not yet rule out a bounded defect from that line. A second count
forces many nonmaximal Ap\'ery elements, which creates the missing
square-root term.

\begin{lemma}\label{lem:paircount}
For any leaf, with $a$ as in \eqref{eq:a},
\begin{equation}\label{eq:paircount}
 \frac{a(a+1)}2\ge2t+a-g.
\end{equation}
\end{lemma}

\begin{proof}
Consider $w\in W$ with $w>F$. Since $S$ is a leaf, $w=u+v$ for positive
$u,v\in S$. The argument used in \cref{lem:leafbasic} works for this $w$
as well: if $u-m\in S$, then $w-m=(u-m)+v\in S$, contradicting $w\in W$.
Thus both summands are in $W$, and both are nonmaximal because each is
strictly below $w$ in the semigroup order.

There are only $a(a+1)/2$ unordered pairs, with repetition, of nonzero
nonmaximal Ap\'ery elements. Their sums can therefore cover at most that
many distinct values of $w$. Hence
\[
 |\{w\in W:w>F\}|\le\frac{a(a+1)}2.
\]
On the other hand, the number of nonzero Ap\'ery elements at most $F$
is at most the number $n=F-g$ of positive semigroup elements at most $F$.
Consequently,
\[
 |\{w\in W:w>F\}|\ge m-1-n.
\]
By $F+t\le2g$ we have $n=F-g\le g-t$, so
\[
 m-1-n\ge t+a-(g-t)=2t+a-g.
\]
Combining these estimates proves \eqref{eq:paircount}.
\end{proof}

\begin{theorem}[Square-root upper bound]\label{thm:upper}
Every leaf of genus $g$ and type $t$ satisfies
\begin{equation}\label{eq:quadratic}
 (2g-3t)^2\ge4t-3,
\end{equation}
and, more explicitly,
\begin{equation}\label{eq:upperexplicit}
 t\le\left\lfloor\frac{6g+2-\sqrt{24g-23}}9\right\rfloor.
\end{equation}
\end{theorem}

\begin{proof}
Set $h=2g-3t$. Rearranging \eqref{eq:paircount} gives
\[
 a^2-a\ge4t-2g=t-h,
 \qquad\text{or}\qquad t\le a^2-a+h.
\]
From \cref{lem:leafbasic}, $1\le a\le(h-1)/2$. The function $x^2-x$ is
increasing on $[1,\infty)$, so
\[
 t\le\left(\frac{h-1}{2}\right)^2-\frac{h-1}{2}+h
   =\frac{h^2+3}{4}.
\]
This proves \eqref{eq:quadratic}.

To obtain the displayed bound on $t$, substitute $t=(2g-h)/3$ to get
\[
 3h^2+4h+9-8g\ge0.
\]
Since $h\ge3$, this implies
\[
 h\ge\frac{\sqrt{24g-23}-2}{3}.
\]
Now use $t=(2g-h)/3$ and the integrality of $t$.
\end{proof}

The decisive count is of \emph{nonmaximal} Ap\'ery elements, not of all
minimal generators. Every high Ap\'ery element needs a decomposition, and
its summands must be nonmaximal Ap\'ery elements. A small supply of such
elements generates only quadratically many sums. Near the extremal line,
there are linearly many high elements to cover, so the supply must grow at
least like a square root. The inequalities above make that observation
quantitative without imposing a special form on $S$.

\section{The maximum type at every large genus}\label{sec:asymptotic}

The unpadded grid family proves a lower bound only at the genera
$g_r=6r^2-3r+4$. To determine an ordinary limit rather than just a limit
superior, one must fill the gaps between these genera. The parameter
$\delta$ in \cref{thm:basis} does exactly that.

\begin{theorem}\label{thm:allg}
For every integer $g\ge7$, there is a leaf of genus $g$. Moreover,
\begin{equation}\label{eq:lowerexplicit}
 L(g)\ge\frac{2g}{3}-\sqrt{24g-87}-2.
\end{equation}
Together with \cref{thm:upper}, this yields
\begin{equation}\label{eq:theta}
 L(g)=\frac23g-\Theta(\sqrt g),\qquad
 \lim_{g\to\infty}\frac{L(g)}g=\frac23.
\end{equation}
\end{theorem}

\begin{proof}
The grid basis with $r=1$ has genus $g_1=7$ and type $3$; arbitrary
$\delta\ge0$ already proves existence at every genus at least $7$.
For the stronger lower bound, choose the largest $r\ge1$ such that
$g_r\le g$, equivalently
\begin{equation}\label{eq:rchoice}
 r=\left\lfloor\frac{3+\sqrt{24g-87}}{12}\right\rfloor.
\end{equation}
Set $\delta=g-g_r$. Since
\[
 g_{r+1}-g_r=12r+3,
\]
we have $0\le\delta\le12r+2$. By \cref{thm:basis}, $S(D_r,\delta)$ has
exactly genus $g$ and still has type $t_r$. Using \eqref{eq:gridratio},
\begin{align*}
 L(g)\ge t_r
 &=\frac23(g-\delta)-4r+\frac73\\
 &\ge\frac23g-12r+1.
\end{align*}
Equation \eqref{eq:rchoice} implies
$12r\le3+\sqrt{24g-87}$, proving \eqref{eq:lowerexplicit}.

Let $E(g)=2g/3-L(g)$. The lower and upper bounds together imply
\begin{equation}\label{eq:errorbounds}
 \frac{\sqrt{24g-23}-2}{9}
 \le E(g)\le\sqrt{24g-87}+2,\qquad g\ge7.
\end{equation}
Both outer expressions have positive square-root leading terms. Hence
$E(g)=\Theta(\sqrt g)$, and division by $g$ proves the asserted limit.
\end{proof}

Here $\Theta(\sqrt g)$ has its usual two-sided meaning: there exist positive
constants $c,C$ and a threshold $g_0$ such that
$c\sqrt g\le2g/3-L(g)\le C\sqrt g$ for every integer $g\ge g_0$.
It does not assert the existence of a limiting square-root coefficient.

There is also an explicit, though not optimized, point beyond which a
counterexample exists at \emph{every} genus.

\begin{corollary}\label{cor:everygenus}
For every integer $g\ge697$, the construction above gives a leaf with
$t>\conjb(g)$.
\end{corollary}

\begin{proof}
Since $g_{11}=697$, the choice in \eqref{eq:rchoice} gives $r\ge11$.
For $0\le\delta\le12r+2$,
\begin{align*}
 2t_r-(g_r+\delta)
 &\ge (2r^2-9r+6)-(12r+2)\\
 &=2r^2-21r+4>0.
\end{align*}
The final polynomial is positive at $r=11$ and increasing thereafter.
Thus $t_r\ge g/2$, and \eqref{eq:beta-half} gives the contradiction.
\end{proof}

The correction to the original picture is therefore not a sporadic
exception or a change of a rounding term. The proposed bound has leading
coefficient $1/2$, whereas the actual extremal function has leading
coefficient $2/3$, with a negative correction of square-root order.

\section{Search, exact verification, and reproducibility}\label{sec:computation}

\subsection{Discovery was optimization; validation is exact}

The initial search used fixed multiplicity $m$ and Frobenius number $F$.
For $1\le i\le F$, binary variables $s_i$ indicated semigroup membership
and binary variables $p_i$ marked proposed pseudo-Frobenius numbers.
Membership was fixed by
\[
 s_1=\cdots=s_{m-1}=0,\qquad s_m=1,\qquad s_F=0.
\]
For indices within $\ints{1}{F}$, closure and the pseudo-Frobenius
implications were encoded by
\begin{align*}
 s_i+s_j-s_{i+j}&\le1&& (i+j\le F),\\
 p_i+s_i&\le1,\\
 p_i+s_j-s_{i+j}&\le1&& (i+j\le F).
\end{align*}
Only $j\ge m$ is needed in the last constraint, and only $i,j\ge m$ in
the closure constraint. Repeated summands have coefficient $2$ where
appropriate.

To ensure the leaf condition, for each $n\in\ints{F+1}{F+m}$ we introduced
binary witnesses $z_{ij}$ for $i+j=n$ with $m\le i\le j\le F$, imposed
$z_{ij}\le s_i,s_j$, and required at least one witness. Such decompositions
cannot use a summand greater than $F$, because then their total with a
positive summand would exceed $F+m$.

The maximized expression was
\[
 2\sum_{i=1}^F p_i+\sum_{i=1}^F s_i-F.
\]
For any fixed membership vector, maximizing the $p_i$ marks every valid
pseudo-Frobenius gap, so this expression becomes $2t-g$. Positive or zero
values are a convenient sufficient test for violating \eqref{eq:oldbound}.
This objective need not find every possible parity-dependent violation.

The optional implementation uses SciPy's mixed-integer linear programming
interface~\cite{SciPy}. An initial run at $(m,F)=(36,71)$ returned a leaf
with $(g,t)=(48,24)$, exceeding the bound $23$. Examining its low elements
led to the restricted-basis construction, and then to the smaller
$(34,17)$ example. The historical candidates are saved in
\texttt{data/discovery\_search.jsonl}.

No floating-point optimality claim is used in the proofs. Search results
are merely proposals: every candidate used in the article is reconstructed
and checked with integer arithmetic. A time-limited optimizer may return a
different candidate, or no candidate, on another machine without affecting
any theorem in this note.

\subsection{Two independent descriptions of the main example}

The principal verifier first starts with the finite-tail description
$S=\{0\}\cup A\cup\tail c$. It checks closure below $c$, finds the exact
gap set, tests the pseudo-Frobenius condition against all positive elements
below $c$, and searches for primitive elements through $F+m$.
Integer bitsets make these tests fast while preserving exactness.

A second method starts only from the listed generators. Form a directed
graph on residues modulo $m$; a generator $a$ gives an edge of weight $a$
from residue $r$ to residue $r+a$. Shortest-path distances from $0$ are
exactly the least representable nonnegative integers in each residue,
namely the Ap\'ery set. The supplied Dijkstra implementation uses integer
weights. Unreachable residues would reveal that the generators do not
have greatest common divisor $1$.

For \eqref{eq:smallS}, the Ap\'ery vector, in residue order modulo $26$, is
\begin{equation}\label{eq:smallapery}
\begin{split}
 (w_0,\ldots,w_{25})={}&(0,27,54,29,56,31,32,59,60,61,62,63,64,\\
 &39,40,41,42,43,44,45,46,47,48,49,50,77).
\end{split}
\end{equation}
Membership is then characterized by
\[
 n\in S\quad\Longleftrightarrow\quad n\ge w_{n\bmod m}
 \qquad(n\ge0).
\]
The maximum entry gives $F=77-26=51$. Counting missing elements residue
by residue gives the independent genus formula
\[
 g=\sum_{i=0}^{m-1}\frac{w_i-i}{m}=34.
\]
Maximal Ap\'ery elements recover the same pseudo-Frobenius set by
\cref{lem:aperyPF}. The code also independently recomputes all primitive
elements from this membership test. Thus the tail description and the
generator description are checked against one another, rather than merely
re-evaluating the same closed formulas twice.

\subsection{Scope of the finite checks}

The default run completed the checks in \cref{tab:checks}; its results are saved
in \path{data/verification_report.json}.

\begin{table}[htbp]
\centering
\begin{tabular}{p{0.68\linewidth}r}
\toprule
Check & Count or range\\
\midrule
Odd-progression and padded square-grid instances & 220\\
Restricted bases with maximum $1\le d\le12$ & 384\\
Those bases checked at $\delta=0,1,7$ & 1,152\\
Distinct examples reconstructed independently by Dijkstra & 26\\
Exhaustive tree levels checked & genus $1$--$20$\\
Semigroups in those levels & 93,141\\
Leaves in those levels & 34,001\\
\bottomrule
\end{tabular}
\caption{Exact finite verification completed for this archive. The infinite
claims are established by the written proofs, not inferred from these counts.}
\label{tab:checks}
\end{table}

For each constructed example, the program checks closure, the leaf condition,
the precise generator set, the genus, the entire pseudo-Frobenius set, the
type formula, and the universal linear and quadratic inequalities.
The exhaustive tree traversal generates children by deleting every primitive
above the Frobenius number. Each child has a unique parent, so starting at
genus one and traversing every child covers every semigroup at the next
genus without duplicates. The original bound holds on all leaves through
genus $20$ in this run.

For orientation, the independently computed small-genus maxima are
\[
\begin{array}{c|rrrrrrrrrrrrrr}
 g&7&8&9&10&11&12&13&14&15&16&17&18&19&20\\\hline
 L(g)&3&3&3&4&5&5&5&6&7&7&7&8&9&9.
\end{array}
\]
These small values explain why limited enumeration does not distinguish the
proposed half-genus scale from a larger eventual scale.

The checks are ordinary executable integer computations, not proofs checked
in a foundational proof assistant. Their purpose is to catch arithmetic,
implementation, and transcription errors, especially in the finite
certificates. They do not certify priority, global minimality of the
counterexample, or optimality of a basis or semigroup beyond the explicitly
enumerated ranges.

\section{What is settled and what remains}\label{sec:remaining}

The original conjecture is settled negatively by \cref{thm:small}. The
same construction supplies infinitely many counterexamples by
\cref{cor:odd}, and counterexamples at every genus at least $697$ by
\cref{cor:everygenus}. The maximum leaf type has the exact leading
coefficient $2/3$, and the magnitude of its deficit from that line is
exactly square-root order by \cref{thm:allg}.

Several more precise questions remain outside those conclusions. First,
the exact values of $L(g)$ beyond the small enumerated range are not
determined here. A constructed type gives a lower bound, not necessarily
the maximum. Second, the ratio
\[
 \frac{2g/3-L(g)}{\sqrt g}
\]
is bounded above and away from zero for large $g$, but the present argument
does not show that it converges. Along the unpadded square-grid sequence,
our examples have deficit $4r-7/3$, whereas padding gives a larger uniform
constant. Better restricted bases, denser choices of their maxima, or
constructions of a different shape could improve the lower-bound constant.

Third, the universal upper bound deliberately forgets information about
which pairs of nonmaximal Ap\'ery elements may have equal sums and which
sums can occur in the relevant residue classes. Retaining that structure
could strengthen the square-root coefficient. The elementary estimate
$a(a+1)/2$ is enough for the order of growth but is not asserted to be
sharp for this problem.

Finally, a genus-minimal counterexample would require either a complete
independent search below genus $34$ or a structural exclusion theorem. The
example here is at the first genus after the source paper's reported range,
but our independent search stops at $20$. No stronger minimality statement
is needed for the disproof or the asymptotic theorem.

The reusable mathematical mechanism is the combination of two finite
objects: a restricted additive basis that guarantees enough decompositions,
and an Ap\'ery poset whose nonmaximal elements account for the cost of those
decompositions. The construction and the upper bound meet at the same
square-root scale, even though they arise from different counts.

\appendix
\section{Additional certificate for the conductor interval}\label{sec:certificate}

To make the generator claim independently inspectable, here are explicit
decompositions of the entire possible-tail-generator interval
$\ints{52}{77}$. The compact proof in \cref{sec:small} groups these into
three blocks; the table lists the individual witnesses.

\begin{center}
\begin{tabular}{r@{\;=\;}l@{\qquad}r@{\;=\;}l}
\toprule
\multicolumn{2}{c}{Lower half} & \multicolumn{2}{c}{Upper half}\\
\midrule
52&$26+26$ &65&$26+39$\\
53&$26+27$ &66&$26+40$\\
54&$27+27$ &67&$26+41$\\
55&$26+29$ &68&$26+42$\\
56&$27+29$ &69&$26+43$\\
57&$26+31$ &70&$26+44$\\
58&$26+32$ &71&$26+45$\\
59&$27+32$ &72&$26+46$\\
60&$29+31$ &73&$26+47$\\
61&$29+32$ &74&$26+48$\\
62&$31+31$ &75&$26+49$\\
63&$31+32$ &76&$26+50$\\
64&$32+32$ &77&$27+50$\\
\bottomrule
\end{tabular}
\end{center}

All summands occur in the proposed generating set. Above $77$, subtracting
$26$ stays in the full tail, completing the argument for every larger
integer. The machine-readable certificate in \path{data/counterexample.json} contains
these decompositions, the witnesses in \cref{tab:witness}, the complete gap
set, the pseudo-Frobenius set, and the residue-ordered Ap\'ery vector.

\section{Reproducing the random draw and the computations}\label{sec:reproduce}

The area-selection method was
\texttt{random.Random(secrets.randbits(128)).randrange(24)}. The sampled seed
was
\[
 \texttt{84421886621299946913970799425001841153}.
\]
The returned zero-based index was $9$, selecting semigroup theory. The
complete ordered list of areas is in \texttt{data/random\_selection.json}.
Replaying the seed verifies the recorded index; it does not require fresh
randomness.

From the extracted archive directory, run
\begin{lstlisting}
python code/verify.py --out data
\end{lstlisting}
The verifier requires only the Python standard library and uses integer
arithmetic. Python 3.10 or later supports its syntax; the saved run used
Python 3.13.5. Verification failures raise exceptions through explicit
checks, so running Python with optimization does not disable them.
The optional argument \texttt{--tree-genus 24} extends the exhaustive
traversal but is not necessary for the default report.

The optional numerical discovery program is separate:
\begin{lstlisting}
python -m pip install -r requirements-search.txt
python code/search.py --multiplicity 36 --frobenius 71 --seconds 15
\end{lstlisting}
The original environment used NumPy 2.3.5 and SciPy 1.17.0. Candidate validity
is checked exactly by the verifier even when a solver stops at its time
limit. Solver optimality is not part of the certificate.

To compile the article, run
\begin{lstlisting}
latexmk -pdf -interaction=nonstopmode -halt-on-error article.tex
\end{lstlisting}
Alternatively, run \texttt{pdflatex article.tex} twice. All references are
included directly in the source, so no bibliography database or external
image is required. The archive includes both the editable source and the
compiled PDF, together with a README, verification data, optional search
requirements, and SHA-256 checksums.

\begin{thebibliography}{9}
\bibitem{CRS}
Jonathan Chappelon, Jorge L. Ram\'irez Alfons\'in, and Dumitru I. Stamate,
\emph{Numerical semigroup tree: type-representation},
arXiv:2507.15006v2, revised 22 November 2025.
\url{https://arxiv.org/abs/2507.15006v2}.
Conjecture 4.4 is the target of this note.

\bibitem{Kohonen}
Jukka Kohonen,
\emph{Early Pruning in the Restricted Postage Stamp Problem},
arXiv:1503.03416, 2015.
\url{https://arxiv.org/abs/1503.03416}.

\bibitem{SciPy}
SciPy developers,
\emph{scipy.optimize.milp --- Mixed-integer linear programming},
SciPy reference documentation, accessed 20 September 2026.
\url{https://docs.scipy.org/doc/scipy/reference/generated/scipy.optimize.milp.html}.
\end{thebibliography}

\end{document}
